\section{Marked permutation modules and changes of coordinates}
Let $r\geq2$, let $V$ be an elementary abelian $2$-group, regarded as an
$\mathbb F_2[S_r]$-module, and let $z\ne0$ be fixed by the stabilizer of
$1$ in $S_r$. Write $s_i=(i\ i+1)$ for $1\leq i<r$. For each $j$ put
$z_j=\pi z$ for any $\pi$ taking $1$ to $j$. This is well-defined:
two such permutations differ on the right by an element fixing $1$.
Let $V_z$ be the span of these $r$ vectors. In $V_z\rtimes S_r$ put
\[
 \rho_0=(z,1),\qquad \rho_i=(0,s_i)\quad(1\leq i<r).
\]
The multiplication convention is $(v,\pi)(w,\sigma)=(v+\pi w,\pi\sigma)$.

\begin{theorem}[A marked module]\label{M-thm:module-marking}
The displayed marks form a string C-group of rank $r$, with generated
group $V_z\rtimes S_r$. This assertion requires no linear independence
of the orbit vectors $z_1,\ldots,z_r$.
\end{theorem}
\begin{proof}
All marks are nonidentity involutions. The marks $\rho_i$ with $i\geq1$
generate the indicated copy of $S_r$; their conjugates of $\rho_0$
generate $V_z$. If $i\geq2$, then $s_i$ fixes $1$ and hence $z$, so
$\rho_0$ commutes with $\rho_i$. The other string commutations are
those of adjacent transpositions.

We prove every subgroup intersection explicitly. For a set of mark
indices $D\subseteq\{0,\ldots,r-1\}$ let $E=D\setminus\{0\}$ and
$P_E=\langle s_i:i\in E\rangle$. This is the product of the symmetric
groups on the connected components of the path with retained edges $E$:
path transpositions generate each component and move no point outside
it. Its orbit on $1$ is an initial segment $[k(E)]$, where $k(E)$ is
one plus the length of the initial run $1,2,\ldots$ contained in $E$.
Set
\[
 W_D=\begin{cases}
 \langle z_1,\ldots,z_{k(E)}\rangle_{\mathbb F_2},&0\in D,\\
 0,&0\notin D.
 \end{cases}
\]
The group generated by $D$ is exactly $W_D\rtimes P_E$. Indeed, all
its vector generators are conjugates of $z$ under $P_E$, and those
conjugates are precisely the vectors in the displayed span. The
other path components fix $z$ and do not introduce any further vector.

For two edge sets $E,F$, a permutation preserving each block of both
component partitions preserves each of their pairwise intersections. Those
intersections are exactly the components of the path retaining
$E\cap F$: two vertices are in the same component precisely when every
edge of their unique interval belongs to both sets. Thus
$P_E\cap P_F=P_{E\cap F}$. Also
$k(E\cap F)=\min\{k(E),k(F)\}$. The vector spaces spanned by initial
segments form a nested chain, even when some orbit vectors are dependent.
Consequently $W_D\cap W_{D'}=W_{D\cap D'}$. Intersecting inside the
Cartesian set $V_z\times S_r$ proves the full C-group property.
\end{proof}

The elementary coordinate change below applies to any group action on
an abelian group, not only to permutation modules.
\begin{lemma}[Cocycle change of coordinates]\label{M-lem:shear}
Let $Q$ act on an abelian group $V$, written additively. Suppose
$c:Q\to V$ satisfies
\[
 c(gh)=c(g)+g c(h).
\]
Then $\Phi_c(v,g)=(v+c(g),g)$ is an automorphism of $V\rtimes Q$,
with inverse $\Phi_{-c}$. It fixes $V$ pointwise and carries any marked
subgroup intersection identity to the corresponding identity on the
image marks. In characteristic two it is an involution.
\end{lemma}
\begin{proof}
The cocycle identity with $g=h=1$ gives $c(1)=0$. For two elements,
\[
 \Phi_c((v,g)(w,h))=(v+gw+c(g)+gc(h),gh)
 =\Phi_c(v,g)\Phi_c(w,h).
\]
Addition of $-c(g)$ in the same coordinate is an inverse. The last
assertions follow from $c(1)=0$, from bijectivity, and from
$c(g)+c(g)=0$ in characteristic two. In particular an automorphism
commutes with taking generated subgroups and their intersections.
\end{proof}

\begin{proposition}[Three pairs of signed markings]\label{M-prop:signed}
Take $V=\mathbb F_2^r$ with its natural permutation action. Write
$e_1,\ldots,e_r$ for its basis and $\lambda=e_1+\cdots+e_r$. For any
of the three seeds
\[
 z=\lambda,\qquad z=e_1,\qquad z=\lambda+e_1,
\]
both of the following lists are string C-group markings:
\[
 (z,1),(0,s_1),\ldots,(0,s_{r-1}),
 \qquad
 (z,1),(\lambda+e_1+e_2,s_1),\ldots,
       (\lambda+e_{r-1}+e_r,s_{r-1}).                 \tag{*}
\]
The common abstract groups for these three pairs are respectively
$C_2\times S_r$, $C_2\wr S_r$, and
\[
 \begin{cases}
 C_2\wr S_r,&r\text{ even},\\
 V^{\mathrm{even}}\rtimes S_r,&r\text{ odd},
 \end{cases}
 \qquad V^{\mathrm{even}}=\{v:\sum_j v_j=0\}.
\]
In the last line $r\geq3$ when odd. The statement includes the odd-rank
second list for the third seed.
\end{proposition}
\begin{proof}
Each seed is nonzero and fixed by the stabilizer of $1$, so the first
list is covered by Theorem~\ref{M-thm:module-marking}. To identify its vector
span, only the third seed needs explanation. Its orbit vectors are
$\lambda+e_j$. A relation with coefficients $b_j$ says
$b=(\sum_j b_j)\lambda$. For even $r$ this forces $b=0$; for odd $r$
the only nonzero relation is $b=\lambda$. In the odd case every orbit
vector has even coordinate sum, so their $(r-1)$-dimensional span is
exactly $V^{\mathrm{even}}$.

Let $a\in V$ have alternating coordinates, so that $a_i+a_{i+1}=1$.
Let $\epsilon:S_r\to\mathbb F_2$ be the sign homomorphism. Define
\[
 c(\pi)=\epsilon(\pi)\lambda+a+\pi a.
\]
The first summand is a cocycle because $\lambda$ is fixed and
$\epsilon$ is additive. The other two are a coboundary:
$a+\pi\sigma a=(a+\pi a)+\pi(a+\sigma a)$.
For $s_i$ the result is $c(s_i)=\lambda+e_i+e_{i+1}$.
Lemma~\ref{M-lem:shear}, in the full ambient $V\rtimes S_r$, therefore
carries the first list in (*) to the second, and fixes its first mark.
It need not preserve the first list's named subgroup as a subset; it
is an isomorphism onto the subgroup generated by the second list.
This distinction matters for odd $r$, where the image copy of $S_r$
may have odd vector coordinates. It does not affect the intersection
property or the asserted abstract group.
\end{proof}

\paragraph{Source correspondence and attribution.}
The ordinary lists in Proposition~\ref{M-prop:signed}, in the order of
its three seeds, are graph (1) of Table 9 and graphs (13), (14) of the
cited Fernandes--Piedade catalogue. The changed-coordinate lists are
graph (2) of Table 9 and graphs (15), (16), respectively. In the bound
arXiv:2407.16003v2 edition the latter four diagrams are in Table 6,
physical page 25; the target paper variously calls that table 5 or 6.
The catalogue's Proposition 4.7 and Table 6 concern $r=n/2\geq7$;
the algebraic marking theorem proved here holds for every $r\geq2$.
Here a vector $e_j$ flips the two points in block $j$, and $s_i$
interchanges adjacent blocks without flips. Thus the vector
$\lambda+e_i+e_{i+1}$ means precisely the vertical $i$-edges on every
block except $i$ and $i+1$, as shown in those diagrams.

The known family statements are Piedade--Tranchida~\cite[Proposition 5.1]{N006}.
The odd-rank graph (16) is expressly left unverified by their method
in the paragraph immediately before that proposition. The formula
above proves that case as well. The source bridge and the full proof were distinctly reviewed in the precursor edition; the present migration retains both. This is not a claim that no previous work settles it. Cocycle
changes of splitting are standard group cohomology; the contribution
here is their explicit use to identify the paired markings, together
with the nested-orbit-span proof which tolerates the odd-rank relation.
No source-specific theorem is presumed in either proof.
